\documentclass[11pt,a4paper]{article}

\usepackage[utf8]{inputenc}
\usepackage[T1]{fontenc}
\usepackage[margin=2.5cm]{geometry}
\usepackage{amsmath,amssymb}
\usepackage{booktabs}
\usepackage{longtable}
\usepackage{array}
\usepackage{xcolor}
\usepackage{enumitem}
\usepackage{listings}
\usepackage[hidelinks]{hyperref}

% Inline code, underscores and similar characters are typeset verbatim
\newcommand{\code}[1]{\texttt{\detokenize{#1}}}

\lstdefinestyle{matlabstyle}{
  language=Matlab,
  basicstyle=\ttfamily\small,
  keywordstyle=\color{blue},
  commentstyle=\color{gray},
  stringstyle=\color{red!60!black},
  showstringspaces=false,
  frame=single,
  breaklines=true,
  columns=fullflexible,
  keepspaces=true
}
\lstset{style=matlabstyle}

\setlist{itemsep=2pt,topsep=4pt}
\setlength{\parindent}{0pt}
\setlength{\parskip}{5pt}

\title{BEADS Visualization of Mixed Quantum States\\[2pt]
       \large MATLAB Code Documentation}
\author{}
\date{}

\begin{document}
\maketitle
\tableofcontents
\newpage

%==========================================================================
\section{Overview}
%==========================================================================
This MATLAB package visualizes \emph{pure and mixed} quantum states of two or three qubits using
the generalized BEADS representation. A mixed state is first decomposed into an ensemble of pure states
whose total entanglement (or connected correlation) is minimal. The BEADS of
the individual pure states are then weighted with their probabilities and
summed. In addition to the bead types of pure states, the code visualizes the genuinely correlation-based
separable structure of the state, its classical and nonclassical parts, and
the total correlation structure.

\paragraph{Only two functions are meant to be used directly.}
Everything else in the code directory is back end.
\begin{itemize}
  \item \code{visualize}: plots the BEADS of one given two- or three-qubit
        density operator (Section~\ref{sec:visualize}).
  \item \code{stacked_plot}: plots the BEADS of a series of two-qubit states
        generated by a parametrized state and stacks them along the
        $z$-axis (Section~\ref{sec:stacked}).
\end{itemize}
Both functions are located in the directory \code{code/visualization}. All
other functions are described in Section~\ref{sec:backend} for completeness
and for developers. A graphical front end for both functions is described in
Section~\ref{sec:app}.

%==========================================================================
\section{Setup}
%==========================================================================
\subsection{Directory layout}
\begin{center}
\begin{tabular}{@{}ll@{}}
\toprule
Directory & Content \\
\midrule
\code{Code/visualization} & the two user functions and all plotting routines \\
\code{Code/helper_functions} & bead calculation and structure analysis (back end) \\
\code{Code/optimizations} & optimization routines (back end) \\
\code{Code/cost_functions} & cost functions of the optimizations (back end) \\
\code{Code/analysis_functions} & entanglement and information measures (back end) \\
\code{Code/execution_files} & scripts which load the spherical harmonics and the Pauli basis \\
\code{Code/code_resources} & data files (operator bases, scalings, color maps) \\
\code{Code/third_party} & third-party code (\code{tight_subplot}) \\
\code{App} & graphical app \code{GeneralizedBEADSApp} (Section~\ref{sec:app}) \\
\code{Documentation} & this document \\
\code{Tests} & example calls (\code{Tests.m}) \\
\bottomrule
\end{tabular}
\end{center}

\subsection{Installation and requirements}
Add the code directory including all subdirectories to the MATLAB path, for
example from the repository root:
\begin{lstlisting}
addpath(genpath("Code"));
\end{lstlisting}
All subdirectories must be on the path, because several functions load the
data files in \code{code_resources} by file name and call the scripts in
\code{execution_files}.

The code uses the following MATLAB features and toolboxes:
\begin{itemize}
  \item MATLAB R2020b or newer: \code{arguments} blocks, the validators
        \code{mustBeTextScalar} and \code{mustBeText}, and
        \code{exportgraphics}.
  \item Optimization Toolbox: \code{fminunc} and \code{fmincon} for the
        optimizations.
  \item Global Optimization Toolbox: \code{GlobalSearch} and
        \code{createOptimProblem} (three-qubit classical beads).
  \item Symbolic Math Toolbox: \code{syms} and \code{solve} for the mapped
        entanglement norms.
  \item Parallel Computing Toolbox: optional. The optimizer options request
        parallel evaluation (\code{UseParallel}).
\end{itemize}

%==========================================================================
\section{Concepts and conventions}
\label{sec:concepts}
%==========================================================================
\subsection{Supported systems and input}
Density operators must be $4\times4$ (two qubits) or $8\times8$ (three
qubits). The function \code{visualize} accepts either a numerical matrix or a
text expression in Dirac notation (Section~\ref{sec:dirac}). In both cases the
input is validated by \code{parse} (square, finite, Hermitian, positive
semidefinite, unit trace, two or three qubits) and the system size is decided
from the size of the resulting density operator.

\subsection{Pauli expansion and structures}
The density operator of $n$ qubits is expanded in the Pauli basis,
\[
  \rho = \frac{1}{2^n}\sum_{P}\langle P\rangle\,P ,
\]
and the terms of the expansion are visualized as spherical functions on the
unit sphere (the beads). Several structures are derived from the Pauli
expansion:
\begin{description}[leftmargin=2.2cm,style=nextline]
  \item[T] total correlation function (all contributions of the expansion).
  \item[E] connected correlation function (entanglement based). Contributions
        which factorize into products of lower-order expectation values are
        removed.
  \item[C] compound correlation function, $C = T - E$.
   \item[M] quantum mutual information (total-correlation-based) structure of the density operator
        where the reducible components of $n$-partite correlation functions
        have been subtracted.
  \item[$\Sigma$ (Sigma)] separable structure of the state: the compound part
        without the contributions that merely reflect products of local
        expectation values, i.e., the genuinely correlation-based separable
        structure.
        
 \item[\texorpdfstring{$M_{\mathrm{Cl}}$}{M\_Cl}] classical part of M,
        obtained by a projective measurement of a target qubit
        with optimized projectors (the projectors for which the conditional entropy of the unmeasured subsystem is minimized).
        
   \item[\texorpdfstring{$E_{\mathrm{Cl}}$}{Sigma\_Cl}] classically accessible part
        of $E$ obtained using the same projectors
    
  \item[\texorpdfstring{$\Sigma_{\mathrm{Cl}}$}{Sigma\_Cl}] classically accessible part
        of $\Sigma$ obtained using the same projectors
        
    \item[D] nonclassical correlation structure, $M - M_{\mathrm{Cl}}$ 
     \item[\texorpdfstring{$E_{\mathrm{D}}$}{E\_D}] nonclassical part
        of $E$, $E - E_{\mathrm{Cl}}$.     
  \item[\texorpdfstring{$\Sigma_{\mathrm{D}}$}{Sigma\_D}] nonclassical part
        of $\Sigma$, $\Sigma - \Sigma_{\mathrm{Cl}}$.
  
\end{description}

\subsection{Beads and their arrangement}
For two qubits the figures show the identity bead (Id), the two single-qubit
beads $Q_1$, $Q_2$ and the bilinear beads of the pair $\{1,2\}$ (an
``even'' bead with the spherical harmonics of ranks 0 and 2 and an ``odd''
bead of rank 1). For three qubits there are, in addition, the single-qubit
bead $Q_3$, the bilinear beads of the pairs $\{1,2\}$, $\{1,3\}$, $\{2,3\}$,
and the trilinear beads $\tau_1,\dots,\tau_4$. Beads whose function values
stay below a tolerance are not drawn.

The radius or the color of a bead encodes the function value. In the default
(non-radial) view all beads are spheres and the function value is encoded by a
red/green color scale (compound part) blended with a yellow/blue color scale
(connected part). In the radial view the distance from the origin encodes the
absolute function value.

\subsection{Connectors}
Connector arcs between the beads $Q_i$ visualize entanglement. The radius of
the arc is proportional to the two-tangle $\tau_{ij}=C_{ij}^2$ (square of the
concurrence), or, optionally in \code{stacked_plot}, to the entanglement of
formation. For three qubits additional connectors run from $Q_1$, $Q_2$, $Q_3$
to $\tau_1$. They encode the three-tangle and the difference between the
tripartite entanglement norm and the three-tangle (the sign of the difference
is encoded in the color of the connectors).

\subsection{Pure state decomposition}
\begin{itemize}
  \item Two qubits: the mixed state is decomposed into pure states by mixing
        its eigenstates with a unitary whose parameters are optimized such
        that the average entanglement (entanglement of formation) of the
        ensemble is minimal. The decomposition can be padded with
        additional zero vectors (\code{PadVectors}) to allow ensembles with
        more states than the rank of the density operator.
  \item Three qubits: the full state is decomposed such that the norm of the
        tripartite connected correlation function is minimal; the
        two-qubit marginals are decomposed with respect to the entanglement
        of formation.
\end{itemize}

\subsection{Target qubits}
The classical bead types ($\Sigma_{\mathrm{Cl}}$, $M_{\mathrm{Cl}}$) and the
nonclassical bead types ($\Sigma_{\mathrm{D}}$, D) depend on the measured
qubit. For two qubits this is one qubit (1 or 2). For three qubits the
vector \code{TargetQubits} has five entries
$[q_{12},q_{13},q_{23},q_1,q_2]$: the target qubit of the marginals $\{1,2\}$,
$\{1,3\}$, $\{2,3\}$ and the two target qubits of the tripartite structure.

\subsection{Bead types}
The bead types (strings) which can be requested are summarized in
Table~\ref{tab:beadtypes}.
\begin{table}[h]
\centering
\begin{tabular}{@{}ll@{}}
\toprule
Bead type & Meaning \\
\midrule
\code{"T"} & total correlation function \\
\code{"E"} & entangling structure correlation function \\
\code{"E_Cl"} & classically accessible part of entangling structure \\
\code{"E_D"} & nonclassical part of entangling structure \\
\code{"C"} & compound correlation function \\
\code{"Sigma"} & separable structure \\
\code{"Sigma_Cl"} & classical part of the separable structure \\
\code{"Sigma_D"} & nonclassical part of the separable structure \\
\code{"M"} & total correlation structure \\
\code{"M_Cl"} & classical correlation structure \\
\code{"D"} & nonclassical correlation structure \\
\bottomrule
\end{tabular}
\caption{Bead types for the option \code{BeadTypes}.}
\label{tab:beadtypes}
\end{table}

%==========================================================================
\section{User function: \texttt{visualize}}
\label{sec:visualize}
%==========================================================================
\subsection{Syntax}
\begin{lstlisting}
visualize(DensityOpExpr, Path)
visualize(DensityOpExpr, Path, "Name", Value, ...)
\end{lstlisting}
\code{DensityOpExpr} and \code{Path} are required. All other inputs are optional
name-value pairs. Names are case-insensitive and can be given as string or
char. The function has no output; the figures are saved to files.

\subsection{Required inputs}
\begin{description}[leftmargin=2.2cm,style=nextline]
  \item[\code{DensityOpExpr}] the density operator, either as a numerical
        matrix ($4\times4$ or $8\times8$) or as a string/char array in Dirac
        notation (Section~\ref{sec:dirac}).
  \item[\code{Path}] the path of the output file. The extension should be
        one of \code{png}, \code{jpg}, \code{pdf}, \code{tiff}. The bead
        type or other suffixes are inserted in front of the extension for
        the individual figures (see below).
\end{description}

\subsection{Text input in Dirac notation}
\label{sec:dirac}
If \code{DensityOpExpr} is a string or char array, it is translated into a
matrix by \code{parse} and then validated. Kets \code{|...>} and bras
\code{<...|} are built from the labels in Table~\ref{tab:labels}, which can be
concatenated for several qubits: \code{|0+R>} is the tensor product of
$|0\rangle$, $|+\rangle$, and $|R\rangle$.
\begin{table}[h]
\centering
\begin{tabular}{@{}ll@{}}
\toprule
Label & State \\
\midrule
\code{0}, \code{1} & computational basis \\
\code{+}, \code{-} & $(|0\rangle\pm|1\rangle)/\sqrt2$ \\
\code{R}, \code{L} & $(|0\rangle\pm i|1\rangle)/\sqrt2$ \\
\code{Phi+}, \code{Phi-} & $(|00\rangle\pm|11\rangle)/\sqrt2$ \\
\code{Psi+}, \code{Psi-} & $(|01\rangle\pm|10\rangle)/\sqrt2$ \\
\code{W} & $(|001\rangle+|010\rangle+|100\rangle)/\sqrt3$ \\
\code{GHZ} & $(|000\rangle+|111\rangle)/\sqrt2$ \\
\bottomrule
\end{tabular}
\caption{Labels in kets and bras (case-sensitive).}
\label{tab:labels}
\end{table}
Besides kets and bras, the expression may contain numbers, \code{i}/\code{j},
\code{pi}, the operators \code{+ - * / ^}, parentheses, transposition
(\code{'}), MATLAB matrix syntax (\code{[1 0; 0 1]}), \code{kron(A,B)}, and the
functions \code{sqrt}, \code{exp}, \code{sin}, \code{cos}, \code{tan},
\code{abs}, \code{conj}, \code{real}, \code{imag}, \code{trace}, \code{diag},
\code{eye}, \code{zeros}, \code{ones}. Other identifiers are rejected. A
combined bracket \code{<a|b>} is read as an inner product.
\begin{itemize}
  \item If the expression evaluates to a column vector (kets only) or a row
        vector (bras only), it is treated as a pure state and cast into its
        projector.
  \item \code{dens(Expr)} builds the projector of a column or row vector
        expression inside a larger expression. Vectors are not normalized
        automatically; an expression with trace different from one is
        rejected.
\end{itemize}
Examples:
\begin{lstlisting}
visualize("dens(|Phi+>)", "bell.png");
visualize("1/2*(|0+R><0+R| + |1L-><1L-|)", "state3.png");
visualize("0.7*dens(|Phi+>) + 0.3*eye(4)/4", "Werner.png");
\end{lstlisting}
Implicit tensor products such as \code{|0>|1>} are not supported; write
\code{|01>} or \code{kron(|0>,|1>)}. Directly adjacent kets and bras, as in
\code{|0><1|}, are multiplied.

\paragraph{Parametrized expressions.}
The function \code{parse} (also used by \code{stacked_plot}) accepts the
name-value input \code{"Parameter"}. A text expression may then contain this
parameter as a variable, and a function handle of the parameter is returned
instead of a matrix, e.g.
\begin{lstlisting}
Handle = parse("p*dens(|Phi+>) + (1-p)*eye(4)/4", "Parameter", "p");
rho = Handle(0.3);
\end{lstlisting}
The parameter name must be a valid MATLAB variable name; \code{R}, \code{L},
and the allowed function names (e.g., \code{kron}, \code{i}, \code{pi},
\code{dens}) are rejected. Since the parameter value is unknown at parse time,
the cast of vectors into projectors and the validation of the density operator
are performed each time the function handle is evaluated.

\subsection{Optional inputs}
\begin{longtable}{@{}p{3.6cm}p{3.6cm}p{7.8cm}@{}}
\toprule
Name & Default & Description \\
\midrule
\endhead
\code{BeadTypes} & all eleven bead types & string or char (array) with the bead
  types to plot (Table~\ref{tab:beadtypes}). \\
\code{PureStateVis} & \code{false} & additionally plot the BEADS of the
  individual pure states of the decomposition. \\
\code{CameraAngle} & \code{[47,30]} (2 qubits), \code{[54,45]} (3 qubits) &
  camera perspective \code{[elevation, rotation]}. Labels are only drawn for
  the default views and for the top view \code{[0,90]}. \\
\code{GridPoints} & \code{[200 400]} & resolution of the spherical grid
  \code{[theta, phi]}. \\
\code{TargetQubits} & \code{[1 1 2 1 2]} & target qubits for the classical
  beads (Section~\ref{sec:concepts}). For two qubits only the first entry is
  used. \\
\code{Id} & \code{false} & draw the identity bead. \\
\code{EntanglementArcs} & \code{true} & draw the connector arcs. \\
\code{RadialPlot} & \code{false} & radial plots, the distance from the origin
  encodes the absolute function value. \\
\code{AutoScale} & \code{false} & scale the function values relative to the
  global maximum; the scale factor is annotated in the figure. \\
\code{AddRadialPlot} & \code{false} & additionally create a radial version of
  each figure (only if \code{RadialPlot} is \code{false}). \\
\code{AutoScaleAddRadial} & \code{false} & autoscale the additional radial
  plots. \\
\code{PadVectors} & $\mathrm{rank}(\rho)^2+1-\mathrm{rank}(\rho)$ & number of
  zero vectors added to the pure state decomposition; the padding is
  switched off for \code{0}. \\
\code{SingleView} & \code{false} & two qubits only: \code{false} creates
  one overview figure of all bead types, \code{true} creates an individual
  figure per bead type. \\
\code{Detailed} & \code{false} & two qubits only: plot the detailed
  overview figure (\code{visualize2_complete_detailed}) instead of the
  standard one. It shows the total correlation beads including the Q-beads
  and the coordinate axes together with the bilinear beads of all bead types.
  Only applies to the overview figure, i.e., if \code{SingleView} is
  \code{false}; ignored otherwise. \\
\code{OptStructure} & \code{0} & integer \code{-1}, \code{0}, or \code{1}:
  structure which the classical beads are optimized for. \code{1} determines
  the degree of freedom of the optimal projector orientations of the
  classical beads by repeated optimizations and uses it to maximize the
  overlap with the separable structure, \code{-1} uses it to minimize the
  overlap (entangling structure), and \code{0} applies no structure
  optimization (the first optimum is used). Applies to two qubits and, for
  three qubits, to the tripartite structure. Considerably increases the
  computation time for \code{-1} and \code{1}. \\
\code{MeasurementDirection} & \code{[]} & measurement direction \code{[x y z]}
  of the classical beads (two qubits and the two-qubit marginals of three
  qubits), used if the optimal direction is arbitrary. If empty, the
  direction is requested in the command window when needed. \\
\code{OptTol} & \code{1e-6} & three qubits only: tolerance (step and function
  tolerance of \code{fminunc}) of the optimization of the pure state
  decomposition (\code{opt_connected_tripartite}). Ignored for two qubits. \\
\code{FontSize} & \code{15} & font size of all labels and annotations. \\
\bottomrule
\end{longtable}

\subsection{Behavior and output files}
\begin{description}[leftmargin=0.5cm]
  \item[Two qubits, \code{SingleView = false}] All nine bead types are
        computed and a single overview figure, which shows the relation of
        the bead types, is saved to \code{Path}. \code{BeadTypes} is
        ignored, and the pure-state figures (\code{PureStateVis}) are not
        created in this mode. If \code{Detailed} is \code{true}, the
        detailed version of the overview figure is created.
  \item[Two qubits, \code{SingleView = true}] One figure per requested bead
        type, saved as \code{<Path>_<BeadType>.<ext>}, plus
        \code{<Path>_<BeadType>_radial.<ext>} if \code{AddRadialPlot} is
        set.
  \item[Three qubits] One figure per requested bead type, saved as
        \code{<Path>_<BeadType>.<ext>} (and the radial variants as above).
  \item[\code{PureStateVis}] The pure states are drawn in groups of subplots
        (two qubits: all states in one figure; three qubits: up to eight
        states per figure), for the bead types T, C, and E which are needed
        for the requested bead types. The files are saved with the suffix
        \code{_pure_decomp<T|C|E>} in the file name.
\end{description}
The computation of the pure state decomposition and of the classical
structures can take considerable time, especially for three qubits.

\subsection{Examples}
Two-qubit Werner state, overview figure:
\begin{lstlisting}
rho = 0.3/4*eye(4) + 0.7*1/2*[1 0 0 1]'*[1 0 0 1];
visualize(rho, "Werner.pdf");
\end{lstlisting}
Two-qubit state, individual figures for selected bead types, with the
radial variant:
\begin{lstlisting}
visualize(rho, "Werner.png", "BeadTypes", ["T","E","C","Sigma_Cl"], ...
          "SingleView", true, "AddRadialPlot", true, "FontSize", 12);
\end{lstlisting}
Three-qubit GHZ state with white noise:
\begin{lstlisting}
ghz = [1;0;0;0;0;0;0;1]/sqrt(2);
rho = 0.8*(ghz*ghz') + 0.2*eye(8)/8;
visualize(rho, "GHZNoise.png", "BeadTypes", ["T","E","M"], ...
          "PureStateVis", true);
\end{lstlisting}

%==========================================================================
\section{User function: \texttt{stacked\_plot}}
\label{sec:stacked}
%==========================================================================
\subsection{Syntax}
\begin{lstlisting}
stacked_plot(ParamState, Path)
stacked_plot(ParamState, Path, "Name", Value, ...)
\end{lstlisting}
\code{stacked_plot} shows how the BEADS of a two-qubit state change along a
one-parameter family of states. The states for a list of parameter values are
stacked on top of each other along the $z$-axis. The labels of the beads are
drawn at the topmost state and the parameter value is written next to each
state. One figure is created for every requested bead type.

\subsection{Required inputs}
\begin{description}[leftmargin=2.2cm,style=nextline]
  \item[\code{ParamState}] a function handle \code{@(z) ...} which returns
        the $4\times4$ density operator for the parameter value \code{z}, or
        a text expression in Dirac notation (Section~\ref{sec:dirac}) which
        contains the parameter named by the option \code{Parameter}, e.g.,
        \code{"p*dens(|Phi+>) + (1-p)*eye(4)/4"}. Text expressions are
        converted by \code{parse}.
  \item[\code{Path}] the path of the output files. The bead type is inserted
        in front of the file extension, i.e., the figures are saved as
        \code{<Path>_<BeadType>.<ext>}.
\end{description}
\pagebreak
\subsection{Optional inputs}
\begin{longtable}{@{}m{3.6cm}p{4cm}p{7.4cm}@{}}
\toprule
Name & Default & Description \\
\midrule
\endhead
\code{ParamValues} & \code{0:0.2:1} & parameter values for which the states
  are computed and plotted. \\
\code{BeadTypes} &
\shortstack[l]{ \code{["M","T","E","C",} \\ \code{"Sigma","Sigma_Cl",} \\
\code{"Sigma_D","D","M_Cl",} \\ \code{"E_D","E_Cl"]}} &
  bead types, one figure per type (Table~\ref{tab:beadtypes}). \\
\code{Id} & \code{false} & draw the identity bead. \\
\code{Parameter} & \code{"z"} & name of the parameter, given without dollar
  signs (they are ignored if present). It is plotted as the LaTeX label
  \code{$name$} and, for a text \code{ParamState}, it is the name of the
  parameter in the expression. \\
\code{FontSize} & \code{12} & font size of all labels and annotations. \\
\code{EntArcs} & \code{true} & draw the connector arcs, scaled by the
  entanglement of each state. \\
\code{EntMeasure} & \code{"tangle"} & entanglement measure for the
  connectors: \code{"tangle"} (two-tangle) or \code{"EoF"} (entanglement of
  formation, in units of bits). \\
\code{GridPoints} & \code{[200, 400]} & resolution of the spherical grid
  \code{[theta, phi]}. \\
\code{OptStructure} & \code{0} & integer \code{-1}, \code{0}, or \code{1}:
  structure which the classical beads are optimized for, \code{1} =
  separable structure (maximal overlap with the separable structure),
  \code{-1} = entangling structure (minimal overlap), \code{0} = no
  structure optimization. \\
\code{MeasurementDirection} & \code{[]} & measurement direction \code{[x y z]}
  of the classical beads, used if the optimal direction is arbitrary. If
  empty, the direction is requested in the command window when needed. \\
\code{TargetQubit} & \code{1} & measured qubit (1 or 2) for the classical
  beads. \\
\bottomrule
\end{longtable}

\subsection{Example}
Werner states with a varying weight $p$ of the Bell state:
\begin{lstlisting}
bell = [1; 0; 0; 1]/sqrt(2);
werner = @(p) (1-p)*eye(4)/4 + p*(bell*bell');
stacked_plot(werner, "WernerStack.png", ...
             "ParamValues", 0:0.25:1, "Parameter", "p", ...
             "BeadTypes", ["T","E","Sigma_Cl"]);
\end{lstlisting}
This creates the files \code{WernerStack_T.png}, \code{WernerStack_E.png},
and \code{WernerStack_Sigma_Cl.png}.

%==========================================================================
\section{Graphical app}
\label{sec:app}
%==========================================================================
The app \code{GeneralizedBEADSApp} (directory \code{App}) is a graphical front
end for \code{visualize} and \code{stacked_plot}. It only collects the inputs,
calls the two functions and displays the figure files which were created; no
file in \code{Code} is modified.

\subsection{Starting the app}
Make the directory \code{App} the current folder (or add it to the path) and
run
\begin{lstlisting}
GeneralizedBEADSApp
\end{lstlisting}
The directory \code{Code} including all subdirectories is added to the path
automatically. The app is implemented as a \code{matlab.apps.AppBase} class
and runs inside MATLAB on all platforms (MATLAB R2020b or newer).

\subsection{Inputs}
\begin{description}[leftmargin=0.5cm]
  \item[Mode] \emph{Single State} calls \code{visualize}, \emph{Parameterized
        State} calls \code{stacked_plot}.
  \item[State] In \emph{Single State}: a density operator as matrix or as
        Dirac-notation expression (Section~\ref{sec:dirac}). In
        \emph{Parameterized State}: a text expression or a function handle
        \code{@(z) ...} of the parameter. The button \emph{Check} validates
        the input.
  \item[Output] Folder (with a browse button), file name and format of the
        figures. The functions insert the suffixes described in
        Section~\ref{sec:visualize}. If files with the resulting names
        already exist, the app asks for confirmation before they are
        overwritten.
  \item[Advanced options] All optional inputs of the function of the current
        mode (the options are listed in Sections~\ref{sec:visualize}
        and~\ref{sec:stacked}). The default value is shown in brackets after
        the label; a field which is left empty or set to \code{(default)} uses
        the default value of the function. Numeric fields accept
        expressions such as \code{[1/sqrt(2) 0 1/sqrt(2)]} or \code{0:0.25:1}.
        The bead types are selected with checkboxes (nothing checked = all).
        The panel can be hidden with the button \emph{Advanced options}.
\end{description}

\subsection{Option labels}
The advanced options of \code{visualize} are shown with the following labels.
\begin{center}
\begin{tabular}{@{}ll@{}}
\toprule
Option & Label in the app \\
\midrule
\code{BeadTypes} & Bead types \\
\code{SingleView} & Single view \\
\code{PureStateVis} & Visualize pure states \\
\code{Detailed} & Plot detailed set of beads \\
\code{CameraAngle} & Camera angle \\
\code{GridPoints} & Grid points \\
\code{TargetQubits} & Target qubits \\
\code{Id} & Show Id-Bead \\
\code{EntanglementArcs} & Entanglement arcs \\
\code{RadialPlot} & Radial plot \\
\code{AutoScale} & Autoscale beads \\
\code{AddRadialPlot} & Additional radial plot \\
\code{AutoScaleAddRadial} & Autoscale beads of radial plot \\
\code{PadVectors} & Added zero vectors for optim. \\
\code{OptStructure} & Maximize structural contribution of M\_Cl-Beads \\
\code{MeasurementDirection} & Default measurement direction \\
\code{OptTol} & Optimization tolerance \\
\code{FontSize} & Font size \\
\bottomrule
\end{tabular}
\end{center}
For \code{stacked_plot}, \code{Parameter} is shown as \emph{Parameter name},
\code{ParamValues} as \emph{Parameter values}, and \code{EntArcs},
\code{EntMeasure}, \code{TargetQubit} as \emph{Entanglement arcs},
\emph{Entanglement measure}, \emph{Target qubit}.

\subsection{Blocked options}
In the mode \emph{Single State}, options without an effect for the current
input are blocked (greyed out) and not passed to \code{visualize}:
\begin{itemize}
  \item \code{BeadTypes}, if the state is a valid two-qubit state and
        \code{SingleView} is \code{false} or default (the overview figure
        contains all bead types).
  \item \code{OptTol}, if the state is a two-qubit state.
  \item \code{Detailed}, if \code{SingleView} is \code{true}.
\end{itemize}
If the state is invalid or empty, the options stay available.

\subsection{Structure optimization}
\emph{Maximize structural contribution of M\_Cl-Beads} is a drop-down menu
with the items \emph{entangling}, \emph{none (default)} and
\emph{separable}, which are passed to \code{OptStructure} as $-1$, $0$ and
$1$, respectively. For \emph{none (default)} nothing is passed. The option is
available in both modes.

\subsection{Measurement direction}
If the field \emph{Default measurement direction} is empty, the app passes the
$z$-axis \code{[0 0 1]} as \code{MeasurementDirection}, so the command-window
prompt of \code{find_projectors_bipartite} never appears.

\subsection{Figure viewer}
After the computation, all figures which were created are loaded into the
viewer. Buttons navigate between the figures, zoom in and out, fit the figure
to the window, and open it externally. The buttons are enabled only if they
can be used (for example, navigation needs at least two figures). PDF files
cannot be previewed.

%==========================================================================
\section{Back end}
\label{sec:backend}
%==========================================================================
The remaining functions are called by the two user functions and are not
meant to be called directly. This section describes their role. The detailed
syntax of each function is given in the header comment of the corresponding
file.

\subsection{Program flow}
\paragraph{\code{visualize}, two qubits.}
\begin{enumerate}
  \item \code{parse} converts and validates the input density operator
        (this is the first step for three qubits, too).
  \item \code{opt_ent_formation} computes the pure state decomposition with
        optimized entanglement of formation.
  \item \code{GetBeads2Qubits} computes all requested beads.
  \item \code{visualize2_complete_view} (overview),\\
        \code{visualize2_complete_detailed} (detailed overview, option
        \code{Detailed}) or\\
        \code{visualize2_universal} (individual bead types) plots the beads;\\
        \code{visualize_purestates2} plots the pure states.
\end{enumerate}
\paragraph{\code{visualize}, three qubits.}
\begin{enumerate}
  \item \code{LoadPaulis} loads the three-qubit Pauli basis.
  \item \code{opt_connected_tripartite} decomposes the full state, and
        \code{opt_ent_formation} the three two-qubit marginals.
  \item \code{GetBeads3Qubits} computes the beads.
  \item \code{visualize3_universal} plots the beads;
        \code{visualize_purestates3} plots the pure states.
\end{enumerate}
\paragraph{\code{stacked_plot}.} For each parameter value,
\code{opt_ent_formation} and \code{GetBeads2Qubits} compute the beads. The
function then draws the stacked figure itself for each bead type.

\subsection{Visualization (\texttt{code/visualization})}
\begin{longtable}{@{}p{5.5cm}p{9.5cm}@{}}
\toprule
Function & Purpose \\
\midrule
\endhead
\code{visualize2_complete_view} & overview figure of all two-qubit bead types
  including the relations between them. \\
\code{visualize2_complete_detailed} & detailed variant of the overview figure
  (option \code{Detailed}): total correlation beads with Q-beads and
  coordinate axes, and the bilinear beads of all bead types in a different
  layout. The figure is cropped at the bottom after plotting. \\
\code{visualize2_universal} & figure of one two-qubit bead type. \\
\code{visualize3_universal} & figure of one three-qubit bead type. \\
\code{visualize_purestates2} & BEADS of all pure states of a two-qubit
  decomposition (weights and entanglement can be annotated). \\
\code{visualize_purestates3} & BEADS of all pure states of a three-qubit
  decomposition (up to eight per figure; two-tangles and three-tangle can be
  annotated). \\
\bottomrule
\end{longtable}

\subsection{Bead calculation (\texttt{code/helper\_functions})}
\begin{longtable}{@{}p{5.5cm}p{9.5cm}@{}}
\toprule
Function & Purpose \\
\midrule
\endhead
\code{GetBeads2Qubits} & spherical function values of all beads of a
  two-qubit state from its pure state decomposition; computes only the beads
  required for the requested bead types. \\
\code{GetBeads3Qubits} & same for three qubits, including the two-qubit
  marginals, the two-tangles, and the three-tangle. \\
\code{traverse_and_multiply} & computes the scaled expansion coefficients of
  an operator in the operator basis (\code{LiSAReal}, \code{LiSAScalings}). \\
\code{rho_connected_pure} & connected and compound part of the density
  operator of a pure state. \\
\code{compound_analysis} & separable structure $\Sigma$ from the compound and
  connected structures. \\
\code{rho_mutualinformation} & mutual-information structure M. \\
\code{split_op_tripartite} & splits a three-qubit operator into its
  bipartite components and the remainder. \\
\code{removeIdLinear} & removes the linear (single-qubit) and optionally the
  identity contributions of an operator. \\
\code{find_projectors_bipartite} & optimized measurement projectors for two
  qubits and the classical and nonclassical structures. \\
\code{find_projectors_tripartite} & optimized measurement projectors for two
  target qubits of three qubits. \\
\code{get_projectors_pref} & projectors for a user-chosen measurement
  direction. \\
\code{get_blended_colors} & color data of a bead from the connected and
  compound function values. \\
\bottomrule
\end{longtable}

\subsection{Mathematical and geometric helpers (\texttt{code/helper\_functions})}
\begin{longtable}{@{}p{5.5cm}p{9.5cm}@{}}
\toprule
Function & Purpose \\
\midrule
\endhead
\code{parse} & converts a numerical matrix or a Dirac-notation expression
  (including \code{dens}) to a density operator and validates it (two or
  three qubits, Hermitian, positive semidefinite, unit trace). \\
\code{partial_trace} & partial trace over one qubit of a two- to four-qubit
  operator. \\
\code{generate_hermitian} & Hermitian matrix from a vector of real
  coefficients (generator of unitary transformations in the optimizations). \\
\code{generate_unitary_haar_1q} & Haar-random single-qubit unitary. \\
\code{extract_euler_angles_1q} & Euler angles (ZYZ convention) of a
  single-qubit unitary. \\
\code{bezier} & cubic Bezier tubes used as connector arcs. \\
\code{get_filled_circle} & surface coordinates of a disk (base of the
  coordinate-axis cones). \\
\code{PermuteCoordinates} & permutes coordinates for drawing the coordinate
  axes. \\
\code{Mat2QuBeads} & converts a matrix to a string in the input format of
  the QuBeads software. \\
\bottomrule
\end{longtable}

\subsection{Optimizations (\texttt{code/optimizations})}
\begin{longtable}{@{}p{5.5cm}p{9.5cm}@{}}
\toprule
Function & Purpose \\
\midrule
\endhead
\code{opt_ent_formation} & pure state decomposition of a two-qubit state
  with minimal average entanglement (entanglement of formation). \\
\code{opt_connected_bipartite} & decomposition with minimal connected
  correlation norm (generic size). \\
\code{opt_connected_tripartite} & decomposition of a three-qubit state with
  minimal tripartite connected correlation norm; the optimization tolerance
  is set by \code{OptTol}. \\
\code{opt_classcorr_bipartite} & measurement projectors on one target qubit
  with maximal classical correlation (minimal conditional entropy). \\
\code{opt_classcorrel3_haar_single} & same for one target qubit of a
  three-qubit state. \\
\code{opt_classcorrel3_haar_seq} & same for two target qubits of a
  three-qubit state, optimized sequentially. \\
\code{opt_classcor_universal3} & optimized projectors for one or two target
  qubits and their application to the structures of the full state and the
  marginals. \\
\code{optPoles} & fits a pair of antipodal poles to a set of Bloch vectors
  (decides whether the optimal projectors are unique). \\
\code{optTip2TipCone} & fits a double cone to a set of Bloch vectors
  (decides whether the optimal projectors have a rotational degree of
  freedom). \\
\bottomrule
\end{longtable}

\subsection{Cost functions (\texttt{code/cost\_functions})}
\begin{longtable}{@{}p{5.5cm}p{9.5cm}@{}}
\toprule
Function & Purpose \\
\midrule
\endhead
\code{cost_ent_formation} & average entanglement of a pure state ensemble
  (cost of \code{opt_ent_formation}). \\
\code{connected_norm_bipartite} & average connected correlation norm of an
  ensemble. \\
\code{connected_norm_tripartite} & the same for three qubits, with
  precomputed Pauli operators. \\
\code{cost_connected} & connected correlation function norm of a two- or
  three-qubit operator. \\
\code{cost_connected_tripartite} & connected correlation function norm of a
  three-qubit operator with precomputed Pauli operators. \\
\code{cost_classcorr} & conditional entropy after a single-qubit
  measurement. \\
\code{cost_classcorr3_haar_seq} & conditional entropy after a single-qubit
  measurement on a three-qubit state (Euler-angle parametrization). \\
\code{cost_find_projectors} & overlap cost for rotating the projectors about
  a fixed axis; the sign is set by \code{OptStructure} (maximal or minimal
  overlap with the separable structure). \\
\code{cost_sigma_cl3} & overlap cost of the rotated projectors with the
  separable structure of a three-qubit state; the sign is set by
  \code{OptStructure}. \\
\bottomrule
\end{longtable}

\subsection{Analysis functions (\texttt{code/analysis\_functions})}
\begin{longtable}{@{}p{5.5cm}p{9.5cm}@{}}
\toprule
Function & Purpose \\
\midrule
\endhead
\code{get_concurrence} & concurrence of a two-qubit state. \\
\code{get_entanglement_of_formation} & entanglement of formation of a
  two-qubit state. \\
\code{get_three_tangle} & two-tangles and three-tangle of a pure three-qubit
  state. \\
\code{get_three_tangle_mixed} & convex-roof three-tangle of a pure state
  ensemble. \\
\code{getBipEntanglementNorm} & mapped bipartite entanglement norm. \\
\code{getEntanglementNorm} & mapped tripartite entanglement norm, sensitive to
  the distinction of GHZ- and W-type entanglement. \\
\code{get_von_neumann_entropy} & von Neumann entropy of a density
  operator. \\
\code{calc_conditionalinfo3} & conditional entropy after a joint measurement
  of two qubits of a three-qubit state. \\
\bottomrule
\end{longtable}

\subsection{Data and scripts}
\begin{itemize}
  \item \code{execution_files/SH_Defs_Real.m}: defines the spherical
        harmonics on the plotting grid. It requires the variables
        \code{res_theta} and \code{res_phi} in the calling workspace.
  \item \code{execution_files/LoadPaulis.m}: loads the three-qubit Pauli
        basis from \code{pauli_basis_dim8.mat} into the workspace.
  \item \code{code_resources}: \code{LiSAReal.mat} and
        \code{LiSAScalings.mat} (operator basis and scalings used for the
        expansion in beads), \code{RGLinearNew20.mat} and
        \code{YBLinearNew20.mat} (color maps), and
        \code{pauli_basis_dim8.mat} (Pauli basis).
  \item \code{third_party/tight_subplot.m}: third-party function for
        subplot layouts.
\end{itemize}

%==========================================================================
\section{Notes and limitations}
%==========================================================================
\begin{itemize}
  \item Only two- and three-qubit systems are supported by
        \code{visualize}; \code{stacked_plot} supports two qubits.
  \item The optimizations use random starting points. Results (in
        particular the orientation of nonunique projectors) can therefore
        vary between runs, and repeated optimizations can be time consuming.
  \item If the optimal measurement direction of the classical beads of a
        two-qubit state (or of a two-qubit marginal) is arbitrary,
        \code{find_projectors_bipartite} uses the option
        \code{MeasurementDirection}. If it is empty, the direction is
        requested in the command window (an empty input selects the
        $z$-axis). The graphical app always passes a direction, the $z$-axis
        if its field is empty, so it never uses the prompt.
  \item The function \code{visualize} inserts suffixes in front of the file
        extension of \code{Path}. Use one of the extensions \code{png},
        \code{jpg}, \code{pdf}, \code{tiff}, otherwise files may be
        overwritten.
  \item Old files which are no longer used by the two user functions are
        located in the \code{ignore} directory outside of \code{code}.
\end{itemize}

\end{document}
